\documentclass[aspectratio=169]{beamer}

\usepackage[T1]{fontenc}
\usepackage[utf8]{inputenc}
\usepackage[brazil]{babel}
\usepackage{lmodern}
\usepackage{graphicx}
\usepackage{xcolor}
\usepackage{hyperref}
\usepackage{listings}

\definecolor{senacblue}{RGB}{0,65,130}
\definecolor{codebg}{RGB}{248,248,248}
\definecolor{codeframe}{RGB}{220,220,220}
\definecolor{codestr}{RGB}{0,128,0}
\definecolor{codekw}{RGB}{0,65,130}
\definecolor{codecom}{RGB}{120,120,120}

\hypersetup{colorlinks=true, urlcolor=senacblue, linkcolor=senacblue}
\pdfstringdefDisableCommands{\def\\{ }}
\newcommand{\code}[1]{\texttt{#1}}

\usetheme{Madrid}
\usecolortheme{dolphin}
\setbeamercolor{structure}{fg=senacblue}
\setbeamercolor{frametitle}{bg=senacblue!5, fg=senacblue}
\setbeamercolor{title}{fg=senacblue}
\setbeamercolor{block title}{bg=senacblue!10, fg=senacblue}
\setbeamercolor{block body}{bg=white}
\setbeamertemplate{navigation symbols}{}

\lstdefinestyle{aula}{
  backgroundcolor=\color{codebg},
  frame=single,
  rulecolor=\color{codeframe},
  basicstyle=\ttfamily\tiny,
  keywordstyle=\bfseries\color{codekw},
  commentstyle=\itshape\color{codecom},
  stringstyle=\color{codestr},
  showstringspaces=false,
  tabsize=2,
  columns=fullflexible,
  breaklines=true,
  breakatwhitespace=true,
  keepspaces=true,
  upquote=true,
  xleftmargin=0.8em,
  framexleftmargin=0.6em,
  framexrightmargin=0.4em
}
\lstset{style=aula}

\title[Fullstack -- AULA 1]{Desenvolvimento Fullstack: Backend com Node.js}
\subtitle{Aula 1 (4h) -- Introdu\c{c}\~ao ao Node.js e ao backend}
\author[Valtemir Proc\'opio]{Valtemir Proc\'opio de Lima\\Senac}
\date[2026]{2026}

\setbeamertemplate{title page}{
  \begin{centering}
    \vspace*{0.6cm}
    \IfFileExists{senac_logo.png}{\includegraphics[height=2.0cm]{senac_logo.png}\par}{}
    \vspace{0.8cm}
    {\usebeamerfont{title}\usebeamercolor[fg]{title}\LARGE\bfseries\inserttitle\par}
    \vspace{0.4cm}
    {\usebeamerfont{subtitle}\large\insertsubtitle\par}
    \vspace{1.0cm}
    {\usebeamerfont{author}\large\insertauthor\par}
    \vspace{0.2cm}
    {\usebeamerfont{date}\normalsize\insertdate\par}
    \vfill
  \end{centering}
}

\setbeamercolor{author in head/foot}{bg=senacblue, fg=white}
\setbeamercolor{title in head/foot}{bg=senacblue!90, fg=white}
\setbeamercolor{date in head/foot}{bg=senacblue, fg=white}

\setbeamertemplate{footline}{%
  \leavevmode%
  \hbox{%
    \begin{beamercolorbox}[wd=.30\paperwidth,ht=2.6ex,dp=1ex,left]{author in head/foot}%
      \hspace*{0.8em}{\color{white}\textbf{\insertshortauthor}}
    \end{beamercolorbox}%
    \begin{beamercolorbox}[wd=.50\paperwidth,ht=2.6ex,dp=1ex,center]{title in head/foot}%
      {\color{white}\textbf{Fullstack -- AULA 1}}
    \end{beamercolorbox}%
    \begin{beamercolorbox}[wd=.20\paperwidth,ht=2.6ex,dp=1ex,right]{date in head/foot}%
      \insertframenumber/\inserttotalframenumber\hspace*{0.8em}
    \end{beamercolorbox}%
  }%
  \vskip0pt%
}

\begin{document}

\begin{frame}[plain]
  \titlepage
\end{frame}

\begin{frame}{Roteiro da Aula}
\tableofcontents
\end{frame}

% ============================================================
\section{O que \'{e} o backend e por que ele existe}
% ============================================================

\begin{frame}{Frontend vs. Backend: a divis\~ao do trabalho}
\begin{columns}[T]
\column{0.48\textwidth}
\begin{block}{Frontend}
\begin{itemize}
  \item Interface visual renderizada no navegador.
  \item HTML, CSS, JavaScript no cliente.
  \item Respons\'avel pela experi\^encia do usu\'ario.
\end{itemize}
\end{block}
\column{0.48\textwidth}
\begin{block}{Backend}
\begin{itemize}
  \item L\'ogica de neg\'ocio, dados e seguran\c{c}a.
  \item Executa no servidor, invis\'ivel ao usu\'ario.
  \item Exposto ao frontend via \textbf{API}.
\end{itemize}
\end{block}
\end{columns}
\vspace{0.4em}
\begin{alertblock}{Regra pr\'atica}
Nunca confie em valida\c{c}\~oes feitas apenas no frontend. O backend \'{e} a \'{u}ltima barreira de seguran\c{c}a.
\end{alertblock}
\end{frame}

\begin{frame}{O fluxo de uma requisi\c{c}\~ao web}
\begin{enumerate}
  \item Usu\'ario interage com a interface (clique, formul\'ario).
  \item Navegador envia uma \textbf{requisi\c{c}\~ao HTTP} ao servidor.
  \item Servidor executa l\'ogica de neg\'ocio (autentica, consulta banco).
  \item Servidor retorna uma \textbf{resposta} (JSON, HTML, arquivo).
  \item Frontend renderiza o resultado ao usu\'ario.
\end{enumerate}
\begin{block}{Protocolo central}
HTTP (HyperText Transfer Protocol) define os verbos, c\'odigos de status e headers que governam toda troca de dados na web.
\end{block}
\end{frame}

\begin{frame}{Verbos HTTP essenciais}
\begin{center}
\begin{tabular}{|c|l|l|}
\hline
\textbf{Verbo} & \textbf{A\c{c}\~ao} & \textbf{Exemplo} \\
\hline
GET    & Buscar/listar recurso   & Listar usu\'arios    \\
POST   & Criar recurso           & Cadastrar usu\'ario  \\
PUT    & Substituir recurso      & Atualizar perfil    \\
PATCH  & Atualizar parcialmente  & Mudar senha         \\
DELETE & Remover recurso         & Excluir conta       \\
\hline
\end{tabular}
\end{center}
\end{frame}

% ============================================================
\section{Node.js: conceitos fundamentais}
% ============================================================

\begin{frame}{O que \'{e} o Node.js?}
\begin{block}{Defini\c{c}\~ao}
Node.js \'{e} um \textbf{runtime} de JavaScript constru\'ido sobre o motor V8 do Chrome, que permite executar JavaScript fora do navegador -- diretamente no servidor.
\end{block}
\begin{itemize}
  \item Criado por Ryan Dahl em 2009.
  \item Baseado em \textbf{I/O n\~ao-bloqueante} e \textbf{event loop}.
  \item Ideal para APIs, microsservi\c{c}os e aplica\c{c}\~oes em tempo real.
  \item Gerenciado pela OpenJS Foundation; vers\~ao LTS recomendada para produ\c{c}\~ao.
\end{itemize}
\end{frame}

\begin{frame}{Por que Node.js para backend?}
\begin{columns}[T]
\column{0.48\textwidth}
\begin{block}{Vantagens}
\begin{itemize}
  \item Mesma linguagem no front e no back.
  \item Alta performance em opera\c{c}\~oes I/O (rede, disco).
  \item Ecossistema gigante via \textbf{npm}.
  \item Comunidade ativa e mercado consolidado.
\end{itemize}
\end{block}
\column{0.48\textwidth}
\begin{block}{Limita\c{c}\~oes a conhecer}
\begin{itemize}
  \item N\~ao \'{e} ideal para processamento CPU-intensivo.
  \item Callback hell (mitigado com async/await).
  \item Erros n\~ao tratados derrubam o processo.
\end{itemize}
\end{block}
\end{columns}
\end{frame}

\begin{frame}{Event Loop: o cora\c{c}\~ao do Node.js}
\begin{block}{Como funciona}
O Node.js possui uma \'{u}nica thread principal. Opera\c{c}\~oes lentas (I/O) s\~ao delegadas a threads auxiliares e o resultado \'{e} retornado via \textbf{callback} quando pronto.
\end{block}
\begin{enumerate}
  \item Requisi\c{c}\~ao entra na \textbf{Call Stack}.
  \item I/O (ex.: leitura de arquivo) \'{e} delegado ao sistema.
  \item A thread principal continua processando outras requisi\c{c}\~oes.
  \item Quando o I/O termina, o callback entra na \textbf{Callback Queue}.
  \item O Event Loop empurra o callback de volta \`a Call Stack.
\end{enumerate}
\end{frame}

% ============================================================
\section{Instala\c{c}\~ao e configura\c{c}\~ao do ambiente}
% ============================================================

\begin{frame}[fragile]{Instalando o Node.js (via NVM -- recomendado)}
\begin{block}{Por que usar NVM?}
O NVM (Node Version Manager) permite instalar e alternar entre m\'ultiplas vers\~oes do Node sem conflitos.
\end{block}
\begin{columns}[T]
\column{0.48\textwidth}
\begin{block}{Linux / macOS}
\begin{lstlisting}
# instalar nvm
curl -o- https://raw.githubusercontent.com/
nvm-sh/nvm/v0.39.7/install.sh | bash

# instalar Node LTS
nvm install --lts
nvm use --lts
node -v   # ex: v22.x.x
npm  -v
\end{lstlisting}
\end{block}
\column{0.48\textwidth}
\begin{block}{Windows}
\begin{lstlisting}
# usar nvm-windows:
# github.com/coreybutler/nvm-windows
nvm install lts
nvm use lts
node -v
npm  -v
\end{lstlisting}
\end{block}
\end{columns}
\end{frame}

\begin{frame}[fragile]{Configurando o VS Code para Node.js}
\begin{itemize}
  \item Extens\~ao \textbf{ESLint} -- verifica erros e padr\~oes de c\'odigo em tempo real.
  \item Extens\~ao \textbf{Prettier} -- formata\c{c}\~ao autom\'atica ao salvar.
  \item Extens\~ao \textbf{REST Client} ou \textbf{Thunder Client} -- testa APIs direto no editor.
  \item Extens\~ao \textbf{GitLens} -- hist\'orico e blame integrado.
\end{itemize}
\begin{block}{settings.json recomendado}
\begin{lstlisting}
{
  "editor.formatOnSave": true,
  "editor.defaultFormatter": "esbenp.prettier-vscode",
  "editor.tabSize": 2
}
\end{lstlisting}
\end{block}
\end{frame}

% ============================================================
\section{Primeiros scripts Node.js}
% ============================================================

\begin{frame}[fragile]{Hello World no Node.js}
\begin{block}{Arquivo: \code{hello.js}}
\begin{lstlisting}[language=JavaScript]
// Primeiro script Node.js
console.log('Hello, Node.js!');

const nome = 'Valtemir';
console.log(`Bem-vindo, ${nome}!`);
\end{lstlisting}
\end{block}
\begin{block}{Executar no terminal}
\begin{lstlisting}
node hello.js
// Hello, Node.js!
// Bem-vindo, Valtemir!
\end{lstlisting}
\end{block}
\end{frame}

\begin{frame}[fragile]{M\'odulos nativos: \code{fs} e \code{path}}
\begin{block}{Lendo um arquivo de texto (\code{fs.readFile})}
\begin{lstlisting}[language=JavaScript]
const fs   = require('fs');
const path = require('path');

const arquivo = path.join(__dirname, 'dados.txt');

fs.readFile(arquivo, 'utf8', (err, conteudo) => {
  if (err) {
    console.error('Erro ao ler arquivo:', err.message);
    return;
  }
  console.log('Conteudo:', conteudo);
});
\end{lstlisting}
\end{block}
\begin{alertblock}{Aten\c{c}\~ao}
O callback \'{e} chamado \textbf{assincronamente} -- o c\'odigo ap\'os \code{readFile} executa antes do arquivo terminar de ser lido.
\end{alertblock}
\end{frame}

\begin{frame}[fragile]{Servidor HTTP nativo (sem framework)}
\begin{block}{Arquivo: \code{server.js}}
\begin{lstlisting}[language=JavaScript]
const http = require('http');

const server = http.createServer((req, res) => {
  res.writeHead(200, { 'Content-Type': 'text/plain; charset=utf-8' });
  res.end('Servidor Node.js funcionando!');
});

server.listen(3000, () => {
  console.log('Servidor ouvindo em http://localhost:3000');
});
\end{lstlisting}
\end{block}
\begin{itemize}
  \item Acesse \url{http://localhost:3000} no navegador para testar.
  \item Use \code{Ctrl+C} para encerrar o servidor.
\end{itemize}
\end{frame}

\begin{frame}[fragile]{Iniciando um projeto com npm}
\begin{block}{Comandos b\'asicos}
\begin{lstlisting}
mkdir meu-backend
cd meu-backend
npm init -y
npm install nodemon --save-dev
npm run dev
\end{lstlisting}
\end{block}
\begin{block}{package.json gerado (trecho)}
\begin{lstlisting}[language=JavaScript]
{
  "name": "meu-backend",
  "scripts": {
    "start": "node server.js",
    "dev":   "nodemon server.js"
  }
}
\end{lstlisting}
\end{block}
\end{frame}

% ============================================================
\section{Sincrono vs. Ass\'incrono: o modelo Node.js}
% ============================================================

\begin{frame}[fragile]{Callback, Promise e Async/Await}
\begin{columns}[T]
\column{0.48\textwidth}
\begin{block}{Callback (estilo antigo)}
\begin{lstlisting}[language=JavaScript]
fs.readFile('a.txt', 'utf8',
  (err, data) => {
    if (err) throw err;
    console.log(data);
  }
);
\end{lstlisting}
\end{block}
\begin{block}{Promise}
\begin{lstlisting}[language=JavaScript]
fs.promises
  .readFile('a.txt', 'utf8')
  .then(data => console.log(data))
  .catch(err => console.error(err));
\end{lstlisting}
\end{block}
\column{0.48\textwidth}
\begin{block}{Async/Await (recomendado)}
\begin{lstlisting}[language=JavaScript]
async function lerArquivo() {
  try {
    const data = await
      fs.promises.readFile(
        'a.txt', 'utf8'
      );
    console.log(data);
  } catch (err) {
    console.error(err);
  }
}
lerArquivo();
\end{lstlisting}
\end{block}
\end{columns}
\end{frame}

\begin{frame}{Por que async/await \'{e} o padr\~ao moderno?}
\begin{itemize}
  \item C\'odigo linear, f\'acil de ler e depurar.
  \item Tratamento de erros com \code{try/catch} familiar.
  \item Funciona em qualquer fun\c{c}\~ao marcada com \code{async}.
  \item Compat\'ivel com todas as APIs que retornam Promise.
\end{itemize}
\begin{alertblock}{Regra de ouro}
Em projetos novos, \textbf{sempre use async/await}. Callbacks ficam restritos a bibliotecas legadas.
\end{alertblock}
\end{frame}

% ============================================================
\section{Estrutura inicial de um projeto backend}
% ============================================================

\begin{frame}[fragile]{Estrutura de pastas recomendada}
\begin{block}{Organiza\c{c}\~ao m\'inima para um projeto Express}
\begin{lstlisting}
meu-backend/
|-- src/
|   |-- routes/       % definicao das rotas
|   |-- controllers/  % logica de cada rota
|   |-- middlewares/  % autenticacao, logs, erros
|   |-- models/       % modelos de dados (BD)
|   `-- app.js        % configuracao do Express
|-- .env              % variaveis de ambiente
|-- .gitignore
|-- package.json
`-- server.js         % ponto de entrada (listen)
\end{lstlisting}
\end{block}
\begin{itemize}
  \item Separar \code{app.js} de \code{server.js} facilita testes automatizados.
  \item \code{.env} nunca deve ser versionado -- adicionar ao \code{.gitignore}.
\end{itemize}
\end{frame}

\begin{frame}[fragile]{Arquivo \code{.gitignore} essencial para Node.js}
\begin{block}{\code{.gitignore}}
\begin{lstlisting}
node_modules/
.env
dist/
*.log
.DS_Store
coverage/
\end{lstlisting}
\end{block}
\begin{alertblock}{Aten\c{c}\~ao cr\'itica}
Jam\'ais versione \code{node\_modules} ou \code{.env}. O \code{node\_modules} tem milhares de arquivos e o \code{.env} pode conter senhas e tokens.
\end{alertblock}
\end{frame}

% ============================================================
\section{Atividade pr\'atica da aula (4h)}
% ============================================================

\begin{frame}{Roteiro da atividade guiada}
\begin{block}{Objetivo}
Configurar o ambiente e criar o primeiro servidor Node.js com rotas diferenciadas por URL.
\end{block}
\begin{enumerate}
  \item \textbf{Instalar} Node.js LTS e configurar VS Code (20 min).
  \item \textbf{Inicializar} projeto com \code{npm init -y} e instalar \code{nodemon} (10 min).
  \item \textbf{Criar} \code{server.js} com servidor HTTP nativo respondendo em 3 rotas (30 min).
  \item \textbf{Refatorar} uso de callbacks para async/await na leitura de arquivo (30 min).
  \item \textbf{Organizar} projeto na estrutura de pastas recomendada (20 min).
  \item \textbf{Publicar} no GitHub com \code{.gitignore} correto (10 min).
\end{enumerate}
\end{frame}

\begin{frame}[fragile]{Desafio: servidor com m\'ultiplas rotas}
\begin{block}{\code{server.js} -- extender com rotas}
\begin{lstlisting}[language=JavaScript]
const http = require('http');

const server = http.createServer((req, res) => {
  res.setHeader('Content-Type',
    'application/json; charset=utf-8');

  if (req.url === '/' && req.method === 'GET') {
    res.writeHead(200);
    res.end(JSON.stringify({ mensagem: 'API online' }));
  } else if (req.url === '/usuarios'
             && req.method === 'GET') {
    res.writeHead(200);
    res.end(JSON.stringify([
      { id: 1, nome: 'Ana' },
      { id: 2, nome: 'Bruno' }
    ]));
  } else {
    res.writeHead(404);
    res.end(JSON.stringify({ erro: 'Rota nao encontrada' }));
  }
});

server.listen(3000, () =>
  console.log('http://localhost:3000'));
\end{lstlisting}
\end{block}
\end{frame}

\begin{frame}{Checkpoint de avalia\c{c}\~ao formativa}
\begin{enumerate}
  \item O servidor responde corretamente nas 3 rotas definidas?
  \item O projeto est\'a organizado na estrutura de pastas recomendada?
  \item O \code{.gitignore} exclui \code{node\_modules} e \code{.env}?
  \item O \code{nodemon} reinicia automaticamente ao editar o arquivo?
  \item O c\'odigo usa \code{async/await} na leitura de arquivo?
  \item O reposit\'orio est\'a publicado no GitHub com commits descritivos?
\end{enumerate}
\end{frame}

\begin{frame}{Entreg\'avel da Aula 1}
\begin{block}{Reposit\'orio GitHub com}
\begin{itemize}
  \item \code{package.json} configurado com scripts \code{start} e \code{dev}.
  \item \code{server.js} com servidor HTTP respondendo em pelo menos 3 rotas.
  \item Estrutura de pastas \code{src/routes}, \code{src/controllers}.
  \item \code{.gitignore} correto (sem \code{node\_modules}).
  \item \code{README.md} descrevendo como rodar o projeto.
\end{itemize}
\end{block}
\begin{alertblock}{Condi\c{c}\~ao para avan\c{c}ar para a Aula 2}
Ambiente funcionando e servidor nativo operacional. A Aula 2 introduz o Express e pressup\~oe este ambiente pronto.
\end{alertblock}
\end{frame}

\begin{frame}{Fechamento da Aula 1}
\begin{block}{S\'intese}
Hoje voc\^e entendeu o papel do backend, instalou o Node.js, compreendeu o Event Loop e criou o primeiro servidor HTTP com rotas diferenciadas.
\end{block}
\begin{itemize}
  \item \textbf{Pr\'oxima aula}: Express -- estrutura, rotas e servidor b\'asico.
  \item \textbf{Tarefa}: publicar o reposit\'orio com o entreg\'avel descrito acima.
\end{itemize}
\end{frame}

% ============================================================
\section{Refer\^encias}
% ============================================================

\begin{frame}[allowframebreaks]{Fontes utilizadas (verificadas)}
\footnotesize
\begin{itemize}
  \item Node.js -- Documenta\c{c}\~ao oficial\\
  \url{https://nodejs.org/en/docs}

  \item npm -- Guia de refer\^encia\\
  \url{https://docs.npmjs.com}

  \item MDN Web Docs -- HTTP overview\\
  \url{https://developer.mozilla.org/pt-BR/docs/Web/HTTP/Overview}

  \item MDN Web Docs -- Introdu\c{c}\~ao a JavaScript ass\'incrono\\
  \url{https://developer.mozilla.org/pt-BR/docs/Learn/JavaScript/Asynchronous/Introducing}

  \item NVM (Node Version Manager)\\
  \url{https://github.com/nvm-sh/nvm}

  \item NVM for Windows\\
  \url{https://github.com/coreybutler/nvm-windows}

  \item OpenJS Foundation -- Node.js release schedule\\
  \url{https://github.com/nodejs/release}

  \item CGI.br -- Cartilha de Seguran\c{c}a para Internet\\
  \url{https://www.cgi.br/publicacao/cartilha-de-seguranca-para-internet/}
\end{itemize}
\end{frame}

\end{document}
